\documentclass[11pt,a4paper]{article}

% ---------- Packages ----------
\usepackage[margin=1in]{geometry}
\usepackage{amsmath, amssymb, amsthm, mathtools}
\usepackage{microtype}
\usepackage{algorithm}
\usepackage{algpseudocode}
\usepackage[colorlinks=true,
            linkcolor=blue!50!black,
            urlcolor=blue!50!black,
            citecolor=blue!50!black]{hyperref}

% ---------- Theorem environments ----------
\theoremstyle{plain}
\newtheorem{theorem}{Theorem}[section]
\newtheorem{proposition}[theorem]{Proposition}
\newtheorem{lemma}[theorem]{Lemma}
\newtheorem{corollary}[theorem]{Corollary}

\theoremstyle{definition}
\newtheorem{definition}[theorem]{Definition}

\theoremstyle{remark}
\newtheorem{remark}[theorem]{Remark}
\newtheorem{example}[theorem]{Example}

% ---------- Math shortcuts ----------
\newcommand{\R}{\mathbb{R}}
\newcommand{\N}{\mathbb{N}}
\newcommand{\E}{\mathbb{E}}
\newcommand{\Var}{\operatorname{Var}}
\newcommand{\Cov}{\operatorname{Cov}}
\newcommand{\Corr}{\operatorname{Corr}}
\newcommand{\F}{\mathcal{F}}
\newcommand{\Prob}{\mathbb{P}}
\newcommand{\Q}{\mathbb{Q}}
\newcommand{\dd}{\mathop{}\!\mathrm{d}}
\newcommand{\eqdef}{\overset{\mathrm{def}}{=}}
\newcommand{\indic}{\mathbf{1}}

% ---------- Document metadata ----------
\title{Monte Carlo Pricing of a European Call:\\
       Control Variates}
\author{Quant Finance Portfolio --- Phase 2, Block 2, Algorithm 2}
\date{\today}

\begin{document}
\maketitle

\begin{abstract}
We apply the control variates technique of the Block~2.0 writeup to
the Monte Carlo pricer of the European call under the exact GBM
sampler of Block~1.1. Two controls are introduced: the discounted
underlying $X_1 = e^{-rT} S_T$ and the discounted asset-or-nothing
payoff $X_2 = e^{-rT} S_T \indic_{\{S_T > K\}}$. Both have closed-form
expectations under GBM. \emph{A naive shape-matching argument suggests
$X_2$ should outperform $X_1$ because it matches the call's payoff
structurally on both the ITM and OTM regions; this prediction is
empirically wrong}. The correct analysis, carried out in detail below,
reveals that $X_2$ carries excess variance unrelated to the target
and ends up with lower Pearson correlation than $X_1$ (typical
values: $\rho_1 \approx 0.92$, $\rho_2 \approx 0.77$). The underlying
control delivers VRF~$\approx 6.9$ at the ATM point; the AON control
delivers only VRF~$\approx 2.5$. The lesson generalises beyond this
example: structural similarity does not imply linear correlation.
\end{abstract}

\tableofcontents

\section{Introduction}
\label{sec:intro}

Block~2.0 derived the control-variate framework: given a target $Y$
with unknown $\E[Y]$ and a control $X$ with known $\E[X]$ correlated
with $Y$, the adjusted estimator
\begin{equation}
\tilde Y(c) \;=\; Y - c\,(X - \E[X])
\end{equation}
is unbiased for $\E[Y]$ for any choice of $c$, with variance
\begin{equation}
\Var(\tilde Y(c^*)) \;=\; \Var(Y)(1 - \rho^2),
\quad
c^* \;=\; \frac{\Cov(Y, X)}{\Var(X)},
\quad
\rho \;=\; \Corr(Y, X).
\end{equation}
The variance reduction factor is $\mathrm{VRF} = 1/(1 - \rho^2)$,
which can be very large when $|\rho|$ is close to $1$.

The challenge in applying CV is finding a control. It must satisfy
two conditions: (a) its expectation must be known in closed form, (b)
it must be highly correlated with the target. The first is
restrictive; the second is what distinguishes a good control from a
mediocre one.

For the European call under GBM, two controls present themselves
naturally: the discounted underlying $X_1 = e^{-rT} S_T$ and the
discounted asset-or-nothing payoff $X_2 = e^{-rT} S_T
\indic_{\{S_T>K\}}$. Both are tractable in closed form. The
expectations of these two controls present a clean opportunity to
illustrate a subtle but important pedagogical point about how
correlation works, which is the main content of this writeup.

\section{Control 1: the discounted underlying}
\label{sec:cv-underlying}

\subsection{Definition and tractability}

The discounted underlying at maturity is
\begin{equation}
X_1 \;\eqdef\; e^{-rT} S_T.
\end{equation}
Under the risk-neutral measure $\Q$, $\{e^{-rt} S_t\}_{t \geq 0}$ is
a martingale, so
\begin{equation}
\E^\Q[X_1] \;=\; e^{-rT} \E^\Q[S_T] \;=\; e^{-rT} \cdot S_0 e^{rT}
\;=\; S_0.
\end{equation}
The expectation is therefore known exactly, with no computation
needed. This is the simplest possible control variate for any
GBM-based pricer.

\subsection{Naive correlation prediction}

The relationship between $Y = e^{-rT} (S_T - K)^+$ and $X_1 = e^{-rT}
S_T$ is:
\begin{itemize}
\item On $\{S_T > K\}$ (ITM): $Y = X_1 - e^{-rT} K$, linear in $X_1$
with slope $1$.
\item On $\{S_T \leq K\}$ (OTM): $Y = 0$, while $X_1 > 0$. Here $Y$
is constant and uncorrelated with $X_1$ locally.
\end{itemize}
A naive prediction based on this dichotomy is that $\rho_1$ should
be moderate, perhaps $0.5$--$0.7$, since one of the two regimes (OTM)
``wastes'' the linear relationship. We will see below that this
prediction \emph{underestimates} $\rho_1$ substantially: the actual
value at the ATM point is $\rho_1 \approx 0.92$.

\section{Control 2: the asset-or-nothing payoff}
\label{sec:cv-aon}

\subsection{Definition and tractability}

The discounted asset-or-nothing call payoff is
\begin{equation}
X_2 \;\eqdef\; e^{-rT}\, S_T \indic_{\{S_T > K\}}.
\end{equation}
Its expectation has a closed Black-Scholes form. From the standard
derivation of the BS call price decomposition,
\begin{equation}
\E^\Q[X_2] \;=\; e^{-rT} \E^\Q\bigl[S_T \indic_{\{S_T > K\}}\bigr]
\;=\; S_0 \, \Phi(d_1),
\end{equation}
where
\begin{equation}
d_1 \;=\; \frac{\log(S_0 / K) + (r + \sigma^2 / 2) T}{\sigma \sqrt T}
\end{equation}
is the standard $d_1$ from Black-Scholes. This identity is the
asset-or-nothing building block of the BS formula
\begin{equation}
C^{\mathrm{BS}} \;=\; S_0 \Phi(d_1) - K e^{-rT} \Phi(d_2),
\quad d_2 = d_1 - \sigma\sqrt T,
\end{equation}
so $\E[X_2]$ is available for free as a byproduct of the BS pricer
code we already implemented in Phase~1.

\subsection{The naive intuition: ``shape matching''}

The relationship between $Y$ and $X_2$ is:
\begin{itemize}
\item On $\{S_T > K\}$: $Y = X_2 - e^{-rT} K$, linear in $X_2$ with
slope $1$.
\item On $\{S_T \leq K\}$: $Y = 0$ \emph{and} $X_2 = 0$
simultaneously, both at the origin.
\end{itemize}
The naive intuition is that $X_2$ should outperform $X_1$ because it
matches the target on \emph{both} regions: zero where the target is
zero, and linearly related where the target is linear. The OTM
``wasted region'' that hurt $X_1$ does not exist for $X_2$; instead
both target and control vanish together. One might therefore predict
$\rho_2$ near $1$, perhaps above $0.95$.

\emph{This prediction is wrong}, and the rest of this writeup is
devoted to understanding why. Structural similarity --- vanishing
together, linear together --- is not the same as global linear
correlation, and the difference is exactly what Pearson's $\rho$
measures.

\section{The empirical surprise and the correct analysis}
\label{sec:correct}

We computed $\rho_1$ and $\rho_2$ analytically (and confirmed
empirically with $10^6$ paths) for the canonical ATM parameters
$(S_0 = K = 100, r = 0.05, \sigma = 0.20, T = 1)$. The results:
\begin{equation}
\rho_1 \;=\; 0.9245, \qquad \rho_2 \;=\; 0.7714,
\end{equation}
\begin{equation}
\mathrm{VRF}_1 \;=\; \frac{1}{1 - \rho_1^2} \;=\; 6.88,
\qquad
\mathrm{VRF}_2 \;=\; \frac{1}{1 - \rho_2^2} \;=\; 2.47.
\end{equation}
The underlying control beats the asset-or-nothing control by almost
a factor of $3$ in variance reduction. The naive prediction was
inverted: $X_1$ is the better control, not $X_2$.

\subsection{The closed-form analysis}

The key algebraic observation is that $Y$ and $X_2$ differ by a
deterministic function of the path:
\begin{equation}
\label{eq:Y-X2}
Y \;=\; X_2 - e^{-rT} K \cdot A,
\qquad A \;\eqdef\; \indic_{\{S_T > K\}}.
\end{equation}
On $\{S_T > K\}$ both sides equal $e^{-rT}(S_T - K)$; on $\{S_T \leq
K\}$ both sides equal zero. So $Y$ is exactly $X_2$ shifted down by a
Bernoulli random variable scaled by the discounted strike.

From this identity:
\begin{equation}
\Var(Y) \;=\; \Var(X_2) - 2 e^{-rT} K \,\Cov(X_2, A) + (e^{-rT} K)^2
\Var(A),
\end{equation}
\begin{equation}
\Cov(Y, X_2) \;=\; \Var(X_2) - e^{-rT} K \,\Cov(X_2, A).
\end{equation}

The closed forms for the second moments under GBM (lognormal MGF
applied to the truncated regions) are:
\begin{align}
\E[X_1] &= S_0, & \Var(X_1) &= S_0^2 \bigl(e^{\sigma^2 T} - 1\bigr), \\
\E[X_2] &= S_0 \Phi(d_1), &
\E[X_2^2] &= S_0^2 \, e^{\sigma^2 T} \, \Phi(d_1 + \sigma\sqrt T),
\end{align}
and the cross terms $\Cov(X_1, A)$, $\Cov(X_2, A)$, $\Cov(X_1, X_2)$
likewise have closed forms in terms of $\Phi(d_1)$, $\Phi(d_2)$, and
$\Phi(d_1 + \sigma \sqrt T)$. We omit the routine algebra and quote
the numerical values for our parameters:
\begin{equation}
\Var(Y) = 216.66, \quad \Var(X_1) = 408.11, \quad \Var(X_2) = 3322.15.
\end{equation}

The most informative comparison is between $\Var(X_1)$ and
$\Var(X_2)$. The variance of the asset-or-nothing control is
\emph{eight times} that of the underlying control, and \emph{fifteen
times} that of the target itself.

\subsection{Why excess variance hurts correlation}

Pearson's correlation is
\begin{equation}
\rho \;=\; \frac{\Cov(Y, X)}{\sqrt{\Var(Y) \cdot \Var(X)}}.
\end{equation}
The covariance and the target variance are fixed by the problem; the
control's variance enters in the denominator. \emph{A control with
excess variance --- variance not explained by linear association
with $Y$ --- gets penalised in the correlation coefficient}.

In our case, the algebraic source of excess variance for $X_2$ is
the multiplicative structure of the indicator $\indic_{\{S_T > K\}}$:
when the call is ITM, the value of $X_2$ depends on $S_T$, which
ranges widely (the typical ITM path can have $S_T$ anywhere from
$K$ to several times $K$); the variance of $X_2$ on the ITM region
is therefore large. The target $Y = X_2 - e^{-rT} K$ on the same
region has the \emph{same} variance as $X_2$ on this region, but the
\emph{global} target variance is much smaller because $Y = 0$ on
OTM (a region of probability $\approx 0.44$). The denominator of
$\rho_2$ therefore picks up the full $\Var(X_2)$ while the numerator
$\Cov(Y, X_2)$ picks up only the portion correlated with $Y$.

For $X_1$, the analogous algebra plays out differently. $X_1$ has
\emph{positive but limited} variance on both ITM and OTM regions:
the variance on OTM is constrained because $S_T \in (0, K]$ there.
This bounded behaviour makes $\Var(X_1)$ small, only $1.9 \times
\Var(Y)$, so the denominator does not penalise the correlation
heavily.

\subsection{The optimal coefficient confirms the picture}

Another way to see the asymmetry is via the optimal coefficient
$c^*$:
\begin{equation}
c_1^* \;=\; \frac{\Cov(Y, X_1)}{\Var(X_1)} \;=\; \frac{274.91}{408.11}
\;=\; 0.674,
\end{equation}
\begin{equation}
c_2^* \;=\; \frac{\Cov(Y, X_2)}{\Var(X_2)} \;=\; \frac{654.44}{3322.15}
\;=\; 0.197.
\end{equation}
For the underlying, $c_1^* \approx 0.67$, indicating that subtracting
$0.67 \times (X_1 - S_0)$ from $Y$ removes about two-thirds of the
``directional'' variance. For the AON, $c_2^* \approx 0.20$, much
smaller: the algorithm ``does not trust'' the AON nearly as much as
it does the underlying, because the AON has so much variance unrelated
to $Y$.

\section{The general lesson}
\label{sec:lesson}

The mistake corrected in Section~\ref{sec:correct} is a subtle one.
The naive intuition --- that $X_2$ matches $Y$ on both regions ---
is correct: $Y$ and $X_2$ are zero together and linear together. But
this kind of \emph{piecewise} matching is not what Pearson's
correlation measures; Pearson's correlation measures \emph{global
linear fit}. Two variables can match each other piecewise and still
have moderate correlation if the piecewise pieces have very
different variances.

A useful diagnostic that would have revealed the problem before
running any code: compute the ratio $\Var(Y) / \Var(X)$ for each
candidate control. A good control has this ratio close to $1$: the
control varies on the same scale as the target. A control with
$\Var(Y)/\Var(X) \ll 1$ has excess variance and will likely give
modest $\rho$.

\begin{center}
\begin{tabular}{l c c c}
Control & $\Var(Y) / \Var(X)$ & Empirical $\rho$ & VRF \\
\hline
$X_1$ (underlying) & $0.53$ & $0.92$ & $6.88$ \\
$X_2$ (AON)        & $0.07$ & $0.77$ & $2.47$ \\
\end{tabular}
\end{center}

The variance-ratio diagnostic correctly predicts the ranking of the
two controls. The naive shape-matching argument does not.

\section{Implementation: empirical \texorpdfstring{$\hat c$}{c-hat}}
\label{sec:empirical}

We use the empirical estimator of $c^*$ from the same paths that
produce the estimator. Given samples $\{(Y_i, X_i)\}_{i=1}^M$:
\begin{equation}
\hat c \;=\; \frac{\sum_{i=1}^M (Y_i - \bar Y)(X_i - \bar X)}
                  {\sum_{i=1}^M (X_i - \bar X)^2}
\quad\text{(OLS slope of $Y$ on $X$)},
\end{equation}
\begin{equation}
\hat C_M^{\mathrm{CV}} \;=\; \frac{1}{M} \sum_{i=1}^M
\bigl(Y_i - \hat c\,(X_i - \E[X])\bigr).
\end{equation}
The bias of $\hat C_M^{\mathrm{CV}}$ is $O(1/M)$, while the standard
error is $O(1/\sqrt M)$, so the bias is asymptotically negligible
for any $M$ at which the standard error is acceptable. See
Block~2.0 writeup, Proposition~4.4.

\subsection{Connection with OLS regression}

The slope estimator $\hat c$ is exactly the slope of the OLS
regression of $Y$ on $X$ with an intercept. The CV-adjusted samples
$\tilde Y_i = Y_i - \hat c (X_i - \E[X])$ are the regression residuals
shifted by $\hat\alpha + \hat c \E[X]$ to recover the original mean.
The variance of $\tilde Y_i$ is the variance of the regression
residuals, which is $\Var(Y)(1 - R^2)$ with $R^2 = \rho^2$ for a
single regressor. The CV variance reduction is therefore exactly
the OLS coefficient of determination.

\section{Algorithmic specification}
\label{sec:algorithm}

\begin{algorithm}[H]
\caption{Control-variate Monte Carlo for the European call (exact GBM)}
\label{alg:mc-cv}
\begin{algorithmic}[1]
\Require Spot $S_0$, strike $K$, maturity $T$, rate $r$, volatility
$\sigma$, sample size $M$, control identifier $c \in \{1, 2\}$,
confidence level $1 - \beta$, random generator \texttt{rng}.
\Ensure Estimate $\hat C^{\mathrm{CV}}$, half-width $h_M^{\mathrm{CV}}$,
sample variance $S_M^{2,\mathrm{CV}}$.
\State $D \gets e^{-rT}$
\For{$i = 1, \ldots, M$}
\State Draw $Z_i \sim \mathcal{N}(0, 1)$ using \texttt{rng}
\State $S_T^{(i)} \gets S_0 \exp((r - \tfrac{1}{2}\sigma^2)T
+ \sigma\sqrt{T} Z_i)$
\State $Y_i \gets D \cdot \max(S_T^{(i)} - K,\, 0)$
\If{$c = 1$ \Comment{control: discounted underlying}}
\State $X_i \gets D \cdot S_T^{(i)}$,\quad $\E X \gets S_0$
\ElsIf{$c = 2$ \Comment{control: asset-or-nothing}}
\State $X_i \gets D \cdot S_T^{(i)} \cdot \indic_{\{S_T^{(i)} > K\}}$
\State $\E X \gets S_0\, \Phi(d_1)$ \Comment{closed BS form}
\EndIf
\EndFor
\State $\bar Y \gets M^{-1} \sum_i Y_i$,\quad $\bar X \gets M^{-1} \sum_i X_i$
\State $\hat c \gets \frac{\sum_i (Y_i - \bar Y)(X_i - \bar X)}
                          {\sum_i (X_i - \bar X)^2}$
\Comment{OLS slope}
\For{$i = 1, \ldots, M$}
\State $\tilde Y_i \gets Y_i - \hat c\,(X_i - \E X)$
\EndFor
\State Reduce: $\hat C^{\mathrm{CV}} \gets M^{-1} \sum_i \tilde Y_i$,
$\quad S_M^{2,\mathrm{CV}} \gets (M-1)^{-1} \sum_i (\tilde Y_i -
\hat C^{\mathrm{CV}})^2$
\State $h_M^{\mathrm{CV}} \gets z_{\beta/2}\,
\sqrt{S_M^{2,\mathrm{CV}}/M}$
\State \Return $(\hat C^{\mathrm{CV}},\, h_M^{\mathrm{CV}},\,
S_M^{2,\mathrm{CV}})$
\end{algorithmic}
\end{algorithm}

\section{Empirical validation strategy}
\label{sec:validation}

The validation script for the CV pricers comprises four tests, run
separately for each control.

\paragraph{Test 1: empirical correlation and VRF.} For each control,
generate $10^6$ paths and measure $\hat\rho = \widehat{\Corr}(Y, X)$
and $\widehat{\mathrm{VRF}} = 1/(1 - \hat\rho^2)$. The
\emph{corrected} predictions, based on the analysis of
Section~\ref{sec:correct}:
\begin{itemize}
\item Control 1 (underlying): $\hat\rho_1 \approx 0.92$, $\mathrm{VRF}_1
\approx 6.9$.
\item Control 2 (AON): $\hat\rho_2 \approx 0.77$, $\mathrm{VRF}_2
\approx 2.5$.
\end{itemize}
The control 1 (underlying) is the better choice for European call
pricing under GBM, contrary to the naive shape-matching prediction.

\paragraph{Test 2: BS coherence.} Each pricer at $M = 50\,000$ paths
should agree with $C^{\mathrm{BS}}$ within a few half-widths.

\paragraph{Test 3: half-width comparison vs IID at fixed budget.}
For each control, run the CV pricer with $M = 100\,000$ paths and
the IID pricer with the same $M = 100\,000$ paths from the same
seed. The squared half-width ratios IID/CV should match the VRFs
from Test~1.

\paragraph{Test 4: input validation.} Mechanical: both pricers
reject bad inputs and accept negative $r$.

\section{Bibliographic notes}
\label{sec:biblio}

The textbook reference is Glasserman~\cite[Section~4.1]{Glasserman}.
The asset-or-nothing decomposition of the BS formula is in
Hull~\cite[Chapter~26]{Hull}. The OLS interpretation of the optimal
control coefficient is well-treated in Asmussen \&
Glynn~\cite[Chapter~V]{AsmussenGlynn}. The pedagogical point about
shape matching versus correlation, while well-known to practitioners,
does not appear to be highlighted in any of the standard textbook
sources we consulted.

\begin{thebibliography}{99}

\bibitem{AsmussenGlynn}
S.\ Asmussen and P.\ W.\ Glynn,
\emph{Stochastic Simulation: Algorithms and Analysis},
Springer, 2007.

\bibitem{Glasserman}
P.\ Glasserman,
\emph{Monte Carlo Methods in Financial Engineering},
Springer, 2004.

\bibitem{Hull}
J.\ C.\ Hull,
\emph{Options, Futures, and Other Derivatives},
10th edition, Pearson, 2018.

\end{thebibliography}

\end{document}
